% =====================================================================
% CAPITULO 3 - DEL PCAP AL DATASET
% =====================================================================
% Capitulo de metodo. Su tesis: dos de los hallazgos centrales del trabajo no
% son de modelado sino de esta fase, y ninguno se habria detectado mirando
% metricas.
%
% ESTILO (31-ago): sin negritas en el texto corrido, sin incisos entre rayas
% salvo cuando la frase los pida, y sin florituras.
% =====================================================================

Los dos capítulos siguientes miden. Este explica sobre qué. Describe cómo se
pasa de una captura de red en bruto a un conjunto de ejemplos etiquetados, y se
detiene especialmente en la parte que suele darse por hecha, que es de dónde
sale la etiqueta.

Esa insistencia no es de trámite. Dos de los hallazgos centrales del trabajo, que
un ataque entero fuera invisible por un problema de extracción y que el tráfico
benigno de referencia estuviera contaminado por el propio atacante, no son
problemas de modelado sino de esta fase. Ninguno se habría detectado mirando
métricas.

\section{Visión general del pipeline}
\label{sec:pipeline}

El proceso tiene cuatro etapas, resumidas en la figura~\ref{fig:pipeline}, y
cada una descarta información de forma irreversible, así que cada descarte
necesita justificación.

\begin{enumerate}
  \item Saneado de la captura. Las capturas del CSE-CIC-IDS2018 contienen
        registros que \texttt{libpcap} rechaza. Se reescribe la captura antes de
        procesarla. La sección~\ref{sec:zeek} explica por qué este paso no es
        opcional.
  \item Extracción con Zeek. Un \emph{script} propio recorre la captura y emite,
        por cada conexión TCP, su carga útil de aplicación agregada junto con los
        metadatos de la conexión.
  \item Vectorización. Cada conexión se convierte en varias representaciones
        numéricas, las vistas que los capítulos siguientes comparan, y se
        almacena en un fichero \texttt{.npz} autocontenido.
  \item Etiquetado. A cada conexión se le asigna la clase correspondiente a
        partir de una \emph{ground-truth} reconstruida sobre las propias
        capturas, no tomada de la documentación del conjunto de datos.
\end{enumerate}

\begin{figure}[htbp]
  \centering
  \includegraphics[width=\textwidth]{F01_pipeline}
  \caption[El pipeline de extracción]{Las cuatro etapas del proceso, desde la
  captura en bruto hasta el conjunto de ejemplos etiquetados.}
  \label{fig:pipeline}
\end{figure}

Dos decisiones atraviesan todo el proceso y se fijan aquí.

\paragraph{La unidad de ejemplo es la conexión TCP.} No el paquete. Cada ejemplo
corresponde a una conexión completa, identificada por el \texttt{uid} que asigna
Zeek. Un ataque se manifiesta a lo largo de una conversación, y tratar cada
paquete como ejemplo independiente haría que un mismo evento apareciera
centenares de veces, repartido además entre entrenamiento y prueba.

\paragraph{El fichero de salida es autocontenido.} Contiene, por flujo, el
histograma de bytes, la secuencia de bytes, la entropía, la 5-tupla de
direcciones y puertos, la marca de tiempo y la etiqueta. Esto permitió más
adelante derivar vistas nuevas, la ráfaga por origen y el inter-arribo, sin
volver a procesar las capturas. Resultó decisivo, porque los datos intermedios
ocupan decenas de gigabytes por día y se borran tras la vectorización.

\section{Extracción del payload con Zeek}
\label{sec:zeek}

Zeek \cite{Zeek} se emplea como analizador de protocolos. Identifica el servicio
por contenido, reensambla el flujo TCP y expone la carga útil de aplicación.
Sobre él se ejecuta un \emph{script} propio que, por cada conexión, acumula
hasta 64 kilobytes de carga útil y los emite junto con el servicio identificado,
la 5-tupla, el número de paquetes, los bytes totales y la duración.

La elección de Zeek frente a alternativas como CICFlowMeter, que es la
herramienta con la que se generan los ficheros de características que acompañan
al conjunto de datos, responde a dos motivos. El primero es que Zeek da acceso a
la carga útil, y no solo a estadísticos agregados, que es el objeto de este
trabajo. El segundo es que su identificación de servicio se basa en el contenido
y no en el puerto, lo que permite la corrección que describe la
sección~\ref{sec:representacion}.

\subsection*{El saneado no es opcional, y descubrirlo costó un ataque entero}

Durante las primeras semanas del trabajo el ataque de fuerza bruta contra SSH no
aparecía en los datos extraídos. El día se procesaba sin errores, Zeek terminaba
correctamente y el resultado contenía únicamente tráfico benigno.

La causa era una truncatura silenciosa. Zeek se detenía a las 18:05 UTC al
encontrar un registro que \texttt{libpcap} rechaza, y descartaba sin avisar todo
lo que venía después, que era precisamente la tarde en que ocurría el ataque. No
se emitía ningún mensaje de error: el proceso terminaba con éxito sobre una
fracción de la captura. Al reescribir la captura antes de procesarla se recuperan
los 2.9 millones de paquetes reales y el ataque aparece completo.

La lección se incorporó al procedimiento como una comprobación obligatoria, que
consiste en validar que la cobertura temporal de lo extraído coincide con el
rango real de la captura. Un fallo así no produce un error, produce un conjunto
de datos plausible y equivocado, que es mucho peor.

\subsection*{Procesamiento por lotes bajo restricción de disco}

Los días completos del CSE-CIC-IDS2018 se distribuyen en centenares de capturas
que suman decenas de gigabytes, y los ficheros intermedios que genera Zeek ocupan
aproximadamente 1.48 veces lo que ocupa la captura de origen. Un día completo no
cabe en disco si se procesa de una vez.

El procesamiento se organiza por lotes. Se ordenan las capturas de mayor a menor,
se procesa cada una y se elimina la captura en cuanto su salida está verificada,
comprobando antes que el número de bytes extraídos es el esperado. El proceso se
detiene de forma limpia si el espacio libre baja de un umbral, y nunca borra nada
si la extracción ha fallado.

Es una restricción de ingeniería, pero condicionó qué días se pudieron procesar y
merece constar. La figura~\ref{fig:volumen} da la escala del trabajo en número de
capturas, gigabytes procesados y flujos resultantes.

\begin{figure}[htbp]
  \centering
  \includegraphics[width=\textwidth]{F10_volumen_procesado}
  \caption[Escala del procesamiento]{Capturas, volumen en disco y flujos
  extraídos por día.}
  \label{fig:volumen}
\end{figure}

\section{Representación vectorial del flujo}
\label{sec:representacion}

De cada conexión se derivan cinco representaciones, que la
figura~\ref{fig:anatomia} resume. Las tres primeras son las que los capítulos
siguientes comparan como vistas; las dos últimas se añadieron al identificar
regímenes que las anteriores no capturaban.

Vale la pena situarlas antes en la taxonomía clásica de la detección de intrusiones,
que distingue tres tipos de características \cite{Lee2000}. Las
\emph{intrínsecas} se derivan de la cabecera de los paquetes de una misma
conexión y la describen por sí sola: duración, número de paquetes y bytes
transportados. Las de \emph{contenido} se derivan de la carga útil, es decir, de
lo que la conexión transporta por encima de la capa de transporte. Las de
\emph{tráfico} no se calculan sobre una conexión sino sobre un grupo de ellas,
agregando por origen o por destino en una ventana temporal.

Las cinco representaciones de este trabajo se reparten sobre esos tres tipos. El
histograma de bytes y la secuencia de bytes son de contenido, y a ellas se añaden
los diecinueve descriptores interpretables que se describen más abajo. Los
metadatos de flujo son intrínsecos. La ráfaga por origen y el inter-arribo son de
tráfico, porque ninguna de las dos se puede calcular mirando una sola conexión.

Hay una colisión de nombres que declarar para que no confunda. Lo que
esta memoria llama metadatos de flujo son características intrínsecas en la
taxonomía anterior, no de tráfico: describen la conexión aislada. El contraste
que atraviesa los capítulos siguientes, entre contenido y flujo, es por tanto un
contraste entre el segundo tipo y el primero, y las características de tráfico
propiamente dichas entran en la vista conductual.

Antes de enumerarlas, una precisión sobre qué se calculan las dos primeras,
porque la unidad es la conexión y no el paquete. El \emph{script} de Zeek emite
una fila por cada segmento TCP que transporte datos, en el orden en que aparece
en la captura y en ambos sentidos. La vectorización recorre esas filas agrupando
por \texttt{uid} y concatena las cargas útiles en ese mismo orden cronológico,
sin separar la dirección cliente-servidor de la inversa. Sobre esa concatenación
se calculan las dos representaciones de contenido: el histograma acumula hasta 64
kilobytes por conexión, que es un tope de memoria, y la secuencia se queda con
los primeros 256 bytes, rellenando con ceros las conexiones más cortas.

Los 256 bytes de la secuencia no son, por tanto, el primer paquete, sino el
comienzo del diálogo. En la práctica suelen abarcar el primer segmento con datos
y parte del siguiente, que es donde vive la cabecera del protocolo de aplicación:
el \emph{banner} SSH, la línea de petición HTTP o el saludo TLS. La
sección~\ref{sec:saliency} muestra que es justamente ahí donde el modelo mira.

\begin{description}
  \item[Histograma de bytes.] Frecuencia de cada uno de los 256 valores posibles
        en la carga útil agregada, normalizada, más la entropía de esa
        distribución. Es una representación sin orden: descarta la secuencia y
        conserva la composición. Sigue la tradición de PAYL \cite{Wang2004}.
  \item[Secuencia de bytes.] Los primeros 256 bytes de la carga útil, en orden,
        como enteros. Conserva la estructura local que el histograma pierde, y es
        la entrada de los modelos convolucionales y recurrentes.
  \item[Metadatos de flujo.] Cuatro escalares de la conexión: número de paquetes,
        bytes totales, longitud de la carga útil y duración. Se les aplica
        $\log(1+x)$ porque sus distribuciones tienen colas muy pesadas.
  \item[Ráfaga por origen.] Número de conexiones que el mismo origen ha abierto
        en el segundo, los cinco segundos y los treinta segundos previos. Se
        calcula de forma causal, mirando solo al pasado, porque un detector en
        línea no puede consultar el futuro.
  \item[Inter-arribo.] Seis características derivadas de los intervalos entre
        conexiones consecutivas del mismo origen: el intervalo anterior, el
        siguiente, la media y el coeficiente de variación en una ventana local,
        una medida de regularidad y la fracción de intervalos cercanos a la
        mediana.
\end{description}

\begin{figure}[htbp]
  \centering
  \includegraphics[width=\textwidth]{F04_anatomia_flujo}
  \caption[De una conexión TCP a los vectores del modelo]{Las representaciones
  que se derivan de una misma conexión.}
  \label{fig:anatomia}
\end{figure}

\subsection*{Descriptores del payload tomados de la literatura}

Al histograma y la secuencia se añade un tercer conjunto de características sobre
la carga útil, esta vez interpretables y tomadas de la literatura de detección
por contenido. Son diecinueve descriptores agrupados en cuatro familias: medidas
de aleatoriedad, como la entropía de Shannon normalizada, el índice de
coincidencia o el ratio de compresión; medidas de forma de la distribución, como
la frecuencia del byte dominante, la masa de los cuatro dominantes o la distancia
a la distribución uniforme; medidas sobre 2-gramas, en la línea de Anagram
\cite{Wang2006}; y medidas léxicas, como la fracción de bytes imprimibles, de
alfanuméricos, de caracteres de control o de bytes no ASCII, orientadas a
detectar inyecciones y \emph{shellcode}.

La motivación es doble. Por un lado permiten unir el análisis de contenido a las
características que produce Zeek en un mismo vector de dimensión manejable. Por
otro son interpretables una a una, cosa que los 256 valores del histograma no
son. El resultado, que el capítulo~\ref{ch:generalizacion} recoge, es que los
diecinueve descriptores rinden aproximadamente igual que las 257 dimensiones del
histograma.

\subsection*{Qué se excluye deliberadamente}

Las direcciones IP y los números de puerto no forman parte de ninguna vista. La
razón se desarrolla en la sección~\ref{sec:fuga}: en este conjunto de datos son
un atajo casi perfecto y además sobreviven a la validación entre días. Se
excluyen desde el principio, salvo en un experimento donde se incluyen
deliberadamente y como control declarado, para medir el tamaño del atajo.

La exclusión no es una precaución conservadora, y se enuncia como lo que
es: usar la dirección de origen como característica no es correcto. No describe
el ataque sino la identidad de quien lo lanzó en este laboratorio concreto, de
modo que un modelo que la use no aprende a detectar nada, aprende una lista de
culpables. Y esa lista no es transferible: cambia con la red, con el escenario y
con el atacante, así que un resultado obtenido con ella no dice nada sobre lo que
el detector haría en otro sitio. Lo mismo vale para el puerto, que depende de la
configuración de cada red. Cuando en esta memoria aparecen direcciones, lo hacen
como identificador para etiquetar y auditar, nunca como entrada del modelo,
excepto en el control positivo declarado de la sección~\ref{sec:por-que-satura} y
en la medición del atajo de la sección~\ref{sec:fuga}, que existen precisamente
para cuantificar el error que se estaría cometiendo.

\subsection*{La selección por servicio, medida}

El servicio se toma del análisis de contenido de Zeek y no del número de puerto.
Esta corrección se aplicó a mitad del trabajo y obligó a reevaluar todos los
resultados anteriores.

La pregunta que quedaba abierta era cuánto cambiaba en la práctica, y se puede
medir: basta comparar, flujo a flujo, la etiqueta de servicio que da el contenido
con la que daría el puerto. La tabla~\ref{tab:servicio-puerto} recoge el
resultado sobre los días en que ambas son comparables.

\begin{table}[htbp]
  \centering
  \caption[Acuerdo entre el servicio identificado por contenido y el puerto]{Para
  cada día, número de flujos en que ambas identificaciones son comparables,
  cuántas coinciden y cuántas se contradicen.}
  \label{tab:servicio-puerto}
  \footnotesize
  \begin{tabular}{lrrrr}
    \toprule
    \textbf{Día} & \textbf{Comparables} & \textbf{Confirman} & \textbf{Acuerdo} & \textbf{Contradicen} \\
    \midrule
    14-02 (SSH)      & 2.200.888 & 2.200.419 & 99.98\,\% &   469 \\
    22-02 (Web)      & 2.769.964 & 2.769.217 & 99.97\,\% &   747 \\
    16-02 (DoS)      & 4.038.401 & 4.037.957 & 99.99\,\% &   444 \\
    04-07 2017 (SSH) &    98.154 &    98.033 & 99.88\,\% &   121 \\
    05-07 2017 (DoS) &   260.796 &   259.410 & 99.47\,\% & 1.386 \\
    06-07 2017 (Web) &    81.352 &    81.330 & 99.97\,\% &    22 \\
    \bottomrule
  \end{tabular}
\end{table}

El acuerdo supera el 99.4\,\% en los seis días, así que ningún resultado
previo cambia. El valor de la corrección es metodológico: los resultados dejan de
depender de una convención de configuración que en otra red podría ser distinta.
La figura~\ref{fig:servicios} muestra, por día, qué protocolos aparecen
identificados por contenido.

\begin{figure}[htbp]
  \centering
  \includegraphics[width=\textwidth]{F07_servicios_por_dia}
  \caption[Servicios presentes en cada día]{Protocolos identificados por
  contenido en cada uno de los días procesados.}
  \label{fig:servicios}
\end{figure}

\section{El problema del etiquetado: el CSV no trae direcciones}
\label{sec:etiquetado}

El CSE-CIC-IDS2018 se distribuye con ficheros CSV de características de flujo ya
etiquetadas. Podría parecer que basta con cruzarlos con las capturas, pero no
basta, por una razón simple: esos CSV no incluyen las direcciones IP. Traen
marcas de tiempo y características agregadas, pero no permiten identificar qué
conexión de la captura corresponde a cada fila.

La etiqueta, por tanto, hay que reconstruirla. El procedimiento adoptado modela
cada ataque como un par $\{\text{atacante}, \text{víctima}\}:\text{puerto}$
acompañado de una ventana temporal, y asigna la clase de ataque a toda conexión
que encaje en ambos criterios. Los ataques y sus parámetros se declaran en un
módulo único, para que la \emph{ground-truth} sea explícita, versionada y
revisable, en lugar de estar dispersa por los \emph{scripts} de evaluación.

El algoritmo~\ref{alg:etiquetado} recoge el procedimiento. Su aparente
sencillez esconde los dos puntos donde este trabajo se equivocó: la comprobación
de la ventana en la línea 4, cuyos límites hay que verificar contra la actividad
real y no contra la documentación, y el hecho de que un flujo que no encaja en
ningún ataque se etiqueta como benigno por defecto, en la línea 9, que es
exactamente lo que contamina la clase negativa cuando la ventana se queda corta.

\begin{algorithm}[htbp]
\caption{Etiquetado de un flujo a partir de la \emph{ground-truth} verificada}
\label{alg:etiquetado}
\begin{algorithmic}[1]
\Require Flujo $f$ con origen $o$, destino $d$, puerto $p$ y marca de tiempo $t$;
         conjunto de ataques $A$ del día
\Ensure Etiqueta de $f$
\ForAll{ataque $a \in A$}
  \State $\mathit{par} \gets \{o, d\}$
  \If{$\mathit{par} \subseteq a.\mathit{direcciones}$ \textbf{y} $p = a.\mathit{puerto}$}
    \If{$a.\mathit{inicio} \le t \le a.\mathit{fin}$}
      \State \Return $a.\mathit{etiqueta}$
    \EndIf
  \EndIf
\EndFor
\State \Return \textsc{Benigno}
\end{algorithmic}
\end{algorithm}

El modelo tiene un límite que declarar, porque determinó qué días se
pudieron usar. Un ataque en el que un único origen contacta con centenares de
destinos distintos no se puede expresar como un par de direcciones, y ese es
exactamente el caso del reconocimiento de la Infiltration. Por eso ese día se
analiza en la sección~\ref{sec:limite} sin construir un conjunto etiquetado.

\subsection*{La zona horaria como fuente de error}

Las capturas están en UTC y la documentación del conjunto de datos en hora local,
con un desfase que no viene declarado de forma inequívoca. Un error de una hora
desplaza la ventana y reetiqueta miles de flujos en ambos sentidos, marcando como
ataque tráfico legítimo y como benigno actividad del atacante.

El desfase se verificó empíricamente en lugar de asumirse. Se localizó el primer
paquete del atacante en la captura y se comparó con la hora declarada. Para el
CSE-CIC-IDS2018 resulta ser UTC$-$4: una anotación de las 08:31
locales corresponde a las 12:30 UTC en la captura. Para el CIC-IDS2017, empleado
en la validación complementaria, resulta ser UTC$-$3. Son distintos, y darlo por
supuesto habría corrompido el etiquetado de tres días.

\section{Ground-truth verificada empíricamente}
\label{sec:ground-truth}

Esta sección recoge la aportación metodológica del capítulo. La
\emph{ground-truth} de los siete días no se tomó de la documentación oficial,
sino que se reconstruyó sobre las capturas y se verificó una por una. La
tabla~\ref{tab:ground-truth} recoge el resultado.

\begin{table}[htbp]
  \centering
  \caption[Ground-truth verificada sobre las capturas]{\emph{Ground-truth}
  verificada empíricamente sobre las capturas, no tomada de la documentación del
  conjunto de datos. Horas en hora local (UTC$-$4).}
  \label{tab:ground-truth}
  \footnotesize
  \begin{tabular}{llll}
    \toprule
    \textbf{Día} & \textbf{Ataque} & \textbf{Atacante} & \textbf{Víctima} \\
    \midrule
    14-02 & SSH-Bruteforce      & \texttt{13.58.98.64}      & \texttt{172.31.69.25:22} \\
    15-02 & GoldenEye/Slowloris & \texttt{18.219.211.138} y \texttt{18.217.165.70} & \texttt{172.31.69.25:80} \\
    16-02 & Hulk/SlowHTTPTest   & \texttt{18.219.193.20}    & \texttt{172.31.69.25:80} \\
    20-02 & DDoS-LOIC-HTTP      & 10 direcciones            & \texttt{172.31.69.25:80} \\
    22-02 & Web (BF/XSS/SQLi)   & \texttt{18.218.115.60}    & \texttt{172.31.69.28:80} \\
    02-03 & Bot (Ares)          & \texttt{18.219.211.138} (C2) & 10 \emph{bots} en \texttt{172.31.69.x} \\
    28-02 & Infiltration        & \texttt{13.58.225.34} (C2)   & \texttt{172.31.69.24} \\
    \bottomrule
  \end{tabular}
\end{table}

\subsection*{Cómo se localiza a un atacante sin conocerlo}

El procedimiento no parte de la dirección documentada, precisamente para poder
contrastarla. Se ordenan los orígenes de la captura por número de conexiones,
porque los ataques de fuerza bruta y volumétricos dominan por varios órdenes de
magnitud, y se contrasta el candidato con la ventana declarada.

Cuando no hay un origen dominante el procedimiento se invierte. Es el caso del
día de la botnet, donde el canal de mando y control mueve poco tráfico y queda
enterrado entre los orígenes legítimos: allí se extrajo primero la carga útil y
se buscó en ella el patrón del canal, para identificar después la dirección. El
día de la Infiltration exigió un tercer procedimiento, sobre el registro de
conexiones y no sobre la carga útil, porque el reconocimiento no exporta datos de
aplicación.

\subsection*{La documentación oficial falló tres veces}

Este es el resultado que justifica el esfuerzo. En tres de los siete días la
documentación del conjunto de datos no coincide con lo que contienen las
capturas.

En el día 14-02 la dirección del atacante SSH documentada no es la que aparece
en la captura, y además la ventana declarada no cubre toda la actividad continua
del atacante.

En el día 16-02 el ataque SlowHTTPTest se dirige al puerto 21 y no al 80, que es
lo que indica la documentación. Es también el ataque que no deja ni un flujo con
carga útil analizable.

En el día 02-03 la documentación parte la actividad de la botnet en dos franjas
horarias. Medida sobre la captura, la actividad es continua, con un hueco máximo
de 32.6 segundos. Etiquetar según la documentación habría marcado como benignas
varias horas de tráfico de la botnet.

No es una crítica al conjunto de datos, que es una contribución valiosa y de uso
generalizado. Es una constatación con consecuencias prácticas: una etiqueta
heredada sin verificar es una fuente de error tan seria como un modelo mal
diseñado, y a diferencia de este último no se manifiesta como un fallo, sino como
un resultado creíble. La literatura reciente ha llegado a la misma conclusión al
reetiquetar CIC-IDS2017 \cite{Engelen2021}.

\subsection*{El desbalanceo que resulta}

La figura~\ref{fig:composicion} muestra la composición de clases de cada día una
vez aplicada la etiqueta verificada. El desbalanceo es extremo y muy desigual
entre días: el día web tiene 203 flujos de ataque entre 2.7 millones, mientras el
del DoS-Hulk tiene 1.8 millones de ataque sobre 4 millones de flujos. Es la razón
por la que el capítulo~\ref{ch:regimenes} balancea 1:1 antes de comparar y no
reporta exactitud global.

\begin{figure}[htbp]
  \centering
  \includegraphics[width=\textwidth]{F06_composicion_clases}
  \caption[Composición de clases por día]{Flujos de ataque frente a flujos
  benignos en cada día, tras aplicar la \emph{ground-truth} verificada.}
  \label{fig:composicion}
\end{figure}

\section{Honestidad del conjunto benigno}
\label{sec:benigno}

Queda un error del que este trabajo fue víctima dos veces, y que merece una
sección propia porque es el más difícil de detectar de todos los descritos.

Cuando la ventana etiquetada no cubre toda la actividad del atacante, las
conexiones que quedan fuera no desaparecen: se quedan en el conjunto de datos
etiquetadas como benignas. El resultado es un conjunto de referencia contaminado
por el propio ataque, y cualquier comparación hecha sobre él enfrenta el ataque
consigo mismo.

La figura~\ref{fig:error-ventana} ilustra el mecanismo, y ocurrió dos veces. En
el día del SSH-Bruteforce, 1.589 de los 2.022 flujos etiquetados como benignos
del puerto 22 eran del propio atacante, y solo 433 procedían de terceros. En el
día del DoS-Hulk, el flujo de ataque es continuo hasta las 17:58:22 UTC mientras
la ventana terminaba a las 17:52:59, con lo que quedaban fuera 740.821 flujos
del atacante, el 65.1\,\% de todo el tráfico HTTP etiquetado como benigno de ese
día.

\begin{figure}[htbp]
  \centering
  \includegraphics[width=\textwidth]{F09_error_ventana}
  \caption[El error de límite de ventana]{Actividad del atacante frente a la
  ventana etiquetada. Las conexiones que quedan fuera no desaparecen: pasan a
  engrosar la clase benigna.}
  \label{fig:error-ventana}
\end{figure}

\paragraph{Por qué no se detecta mirando las métricas.} Porque los resultados que
produce son plausibles. Con el benigno contaminado, ninguna vista superaba el
$0.78$ en el día del SSH, y ese techo se interpretó durante semanas como una
propiedad del problema, como si los metadatos no dieran para más, cuando era un
artefacto del etiquetado. Un número bajo no levanta sospechas del mismo modo que
un $1.0000$, y ahí está la trampa.

\paragraph{Cómo se detecta.} Auditando el origen de cada flujo de la clase
negativa, y comprobando que la ventana etiquetada cubre la actividad continua del
atacante sobre el conjunto completo y no sobre una captura suelta. La
figura~\ref{fig:composicion-benigno} descompone el tráfico etiquetado como
benigno según su procedencia real. La corrección consistió en arreglar las
ventanas y reetiquetar esos flujos como lo que eran, actividad del ataque.

\begin{figure}[htbp]
  \centering
  \includegraphics[width=\textwidth]{F15_composicion_benigno}
  \caption[De qué estaba hecho el tráfico benigno]{Descomposición del tráfico
  etiquetado como benigno según su procedencia real.}
  \label{fig:composicion-benigno}
\end{figure}

\paragraph{Una consecuencia para la reproducibilidad.} Al corregir el etiquetado
en el propio conjunto de datos, el escenario contaminado deja de ser
reconstruible. Las cifras que en su momento lo describieron son históricas y no
comparables con las de los capítulos siguientes, y así se declaran cada vez que
se mencionan.

Esta es, junto con la verificación de la \emph{ground-truth}, la razón por la que
este capítulo ocupa el lugar que ocupa en la memoria. Los dos capítulos que
siguen no tendrían valor sin él, porque medirían con precisión sobre datos
equivocados.
